\documentclass{article}
\usepackage[T1]{fontenc}
\usepackage{lmodern}
\usepackage{graphicx}
\usepackage{amsmath,amssymb}
\usepackage[a4paper,margin=2cm]{geometry}
\usepackage[absolute,overlay]{textpos}
\setlength{\TPHorizModule}{1mm}
\setlength{\TPVertModule}{1mm}
\usepackage[indonesian]{babel}
\usepackage{hyperref}
\setlength{\emergencystretch}{2em}
% unit: Halaman 100

\begin{document}

% === Halaman 100 (python/native) ===

100

K–1  =













9

15

10

3

7

1

 = 













−

−

−

3

15

10

9

7 1


                                    = 













=













=













=













−

−

19

9

6

5

26

mod

45

165

240

135

3

11

16

9


15

3

15

10

9

15


Keterangan (ingat kembali teori bilangan di dalam Matematika Diskrit):

 (i) 7 –1 (mod 26)  15, karena (7)(15) = 105 mod 26 = 1

 (ii) –10  16 (mod 26)

 (iii) –15  11 (mod 26)    )

\end{document}
